\documentclass[12pt]{article}

\usepackage{sbc-template}
\usepackage{graphicx,url}
\usepackage[brazil]{babel}
\usepackage[utf8]{inputenc}
\usepackage[T1]{fontenc}
\usepackage{enumitem}
\usepackage{placeins}
\usepackage{booktabs}

\sloppy

\title{Constru\c{c}\~ao Semi-Autom\'atica de um L\'exico de Misoginia em
       Portugu\^es Brasileiro: Resultados Preliminares}

\author{Hugo Guilherme de Assis Paula\inst{1}
  , Deborah Silva Alves Fernandes\inst{1}}

\address{Instituto de Inform\'atica -- Universidade Federal de Goi\'as (UFG)\\
  Goi\^ania -- GO -- Brasil
  \email{hugo.guilherme.paula@gmail.com, deborah@inf.ufg.br}
}

\begin{document}

\maketitle

\begin{resumo}
A detec\c{c}\~ao autom\'atica de misoginia em ambientes digitais demanda
recursos lexicais especializados em portugu\^es brasileiro, ainda escassos
na literatura. Este artigo descreve a constru\c{c}\~ao preliminar de um
l\'exico de termos associados a textos com tra\c{c}os de misoginia, induzido
de forma semi-autom\'atica a partir de um corpus binariamente anotado de
74.338 documentos. O m\'etodo combina sele\c{c}\~ao por TF--IDF
discriminativo e expans\~ao por PMI, com normaliza\c{c}\~ao via
lematiza\c{c}\~ao do spaCy (\texttt{pt\_core\_news\_lg}). O recurso
preliminar cont\'em 531 termos. Em uma avalia\c{c}\~ao bin\'aria
out-of-sample, um classificador linear com features derivadas do l\'exico
atinge F1 = 0,87 na classe de interesse, indicando que o recurso \'e
promissor como sinal complementar para tarefas de detec\c{c}\~ao
autom\'atica.
\end{resumo}

\begin{abstract}
The automatic detection of misogyny in digital environments requires
specialized lexical resources in Brazilian Portuguese, still scarce in the
literature. This paper describes the preliminary construction of a lexicon
of terms that may be associated with texts displaying traces of misogyny,
induced semi-automatically from a binary-annotated corpus of 74,338
documents. The method combines discriminative TF--IDF selection and PMI
expansion, with normalization via spaCy lemmatization
(\texttt{pt\_core\_news\_lg}). The preliminary resource contains 531 terms.
In an out-of-sample binary evaluation, a linear classifier with
lexicon-derived features reaches F1 = 0.87 on the target class, indicating
the resource is promising as a complementary signal for automatic detection
tasks.
\end{abstract}

% ===========================================================================
\section{Introdu\c{c}\~ao}
% ===========================================================================

A misoginia em ambientes digitais constitui um problema relevante para a
an\'alise autom\'atica de linguagem, especialmente quando se busca
distinguir ofensa gen\'erica de ataques direcionados a mulheres. No
contexto brasileiro, esse fen\^omeno se manifesta em redes sociais e
plataformas de coment\'ario com vocabul\'ario espec\'ifico, frequentemente
ancorado em estere\'otipos culturais e na sobreposi\c{c}\~ao com o discurso
pol\'itico de \'odio, o que dificulta a captura por listas gen\'ericas de
termos ofensivos.

Apesar do avan\c{c}o das abordagens de classifica\c{c}\~ao para discurso
de \'odio e linguagem abusiva em portugu\^es brasileiro
\cite{vargas2022hatebr,leite2020toldbr}, ainda h\'a escassez de
\emph{recursos lexicais} especializados em misoginia. Tal lacuna limita
aplica\c{c}\~oes de PLN que dependem de representa\c{c}\~ao lingu\'istica
interpret\'avel, como classificadores baseados em \emph{features}
estruturadas ou m\'odulos de explicabilidade.

Antes de prosseguir, adotamos a seguinte \textbf{defini\c{c}\~ao
operacional}: neste artigo, consideram-se \emph{textos com tra\c{c}os de
misoginia} aqueles que apresentam hostilidade, desqualifica\c{c}\~ao,
sexualiza\c{c}\~ao ofensiva ou agress\~ao verbal direcionada a mulheres;
os demais textos s\~ao tratados como n\~ao-mis\'oginos para fins da tarefa
bin\'aria adotada. A express\~ao ``tra\c{c}os de'' refor\c{c}a que a
classifica\c{c}\~ao final depende da ocorr\^encia em contexto: palavras
isoladas n\~ao s\~ao mis\'oginas por si, mas podem indicar a presen\c{c}a
de um discurso mis\'ogino quando observadas em texto real.

Diante dessa lacuna, este trabalho tem como objetivo construir
preliminarmente um l\'exico de misoginia em portugu\^es brasileiro por
meio de indu\c{c}\~ao semi-autom\'atica a partir de corpora anotados e
avaliar seu uso em uma tarefa bin\'aria de classifica\c{c}\~ao textual.
As contribui\c{c}\~oes deste artigo s\~ao:

\begin{enumerate}[leftmargin=*,topsep=2pt,itemsep=2pt]
  \item Um \textbf{pipeline preliminar de indu\c{c}\~ao lexical} que
        combina sele\c{c}\~ao por TF--IDF discriminativo e expans\~ao por
        PMI sobre lemas, descrito de forma reprodut\'ivel.
  \item Um \textbf{l\'exico inicial} de 531 termos em portugu\^es
        brasileiro, organizados em classes lexicais sugeridas com base em
        taxonomia consolidada.
  \item Uma \textbf{avalia\c{c}\~ao inicial} do l\'exico em
        classifica\c{c}\~ao bin\'aria \emph{out-of-sample}, comparando
        modelos de refer\^encia com a configura\c{c}\~ao apoiada pelo
        recurso lexical.
\end{enumerate}

% ===========================================================================
\section{Trabalhos Relacionados}
% ===========================================================================

Trabalhos anteriores abordam a detec\c{c}\~ao autom\'atica de misoginia
e de linguagem abusiva por meio de modelos supervisionados, recursos
lexicais e corpora anotados em diferentes idiomas e contextos digitais.

No \^ambito da \textbf{detec\c{c}\~ao de misoginia}, Anzovino, Fersini e
Rosso \cite{anzovino2018} prop\~oem a primeira taxonomia operacional
voltada \`a classifica\c{c}\~ao autom\'atica em \emph{tweets}, com cinco
classes (\emph{discredit}, \emph{stereotype \& objectification},
\emph{sexual harassment \& threats of violence}, \emph{dominance} e
\emph{derailing}). Sultana, Sarker e Bosu \cite{sultana2021} estendem esse
quadro com uma r\'ubrica de anota\c{c}\~ao que separa classes adicionais
de hostilidade impl\'icita (\emph{victim blaming}, \emph{objectification})
e operacionaliza crit\'erios verific\'aveis por anotadores humanos.

Em \textbf{recursos para portugu\^es brasileiro}, dois corpora s\~ao
particularmente relevantes para nosso trabalho. O HateBR
\cite{vargas2022hatebr} re\'une 7.000 coment\'arios anotados do
Instagram, com r\'otulos hier\'arquicos de linguagem ofensiva e tipos de
\'odio, incluindo misoginia. O ToLD-BR \cite{leite2020toldbr} contribui
com 21.000 \emph{tweets} anotados separadamente para cada tipo de
toxicidade (misoginia, racismo, homofobia, insulto, obsceno, xenofobia),
o que permite isolar a dimens\~ao misoginia de outras formas de
hostilidade.

No \textbf{m\'etodo de constru\c{c}\~ao lexical}, Monteles
\cite{monteles2023} prop\~oe um pipeline de expans\~ao autom\'atica de
l\'exico para an\'alise de sentimentos em \emph{tweets} que combina uma
\emph{semente} obtida por TF--IDF discriminativo com uma fase de
expans\~ao baseada em PMI. Adotamos esse m\'etodo, com a substitui\c{c}\~ao
de \emph{stemming} por \textbf{lematiza\c{c}\~ao do spaCy}
\cite{honnibal2017spacy} para garantir legibilidade dos termos --
crit\'erio importante porque o recurso ser\'a submetido a valida\c{c}\~ao
manual por anotadores humanos.

Diferentemente de abordagens centradas exclusivamente em classifica\c{c}\~ao
supervisionada de ponta a ponta, este estudo concentra-se na
\textbf{constru\c{c}\~ao de um recurso lexical especializado} em
portugu\^es brasileiro, com avalia\c{c}\~ao inicial de seu potencial
discriminativo.

% ===========================================================================
\section{Metodologia}
% ===========================================================================

A metodologia foi organizada em quatro etapas: defini\c{c}\~ao operacional
do fen\^omeno, prepara\c{c}\~ao do corpus, constru\c{c}\~ao do l\'exico
e avalia\c{c}\~ao preliminar em classifica\c{c}\~ao bin\'aria.

\subsection{Defini\c{c}\~ao operacional}

Conforme apresentado na Se\c{c}\~ao 1, consideram-se textos com tra\c{c}os
de misoginia aqueles que apresentam hostilidade, desqualifica\c{c}\~ao,
sexualiza\c{c}\~ao ofensiva ou agress\~ao verbal direcionada a mulheres.
A taxonomia adotada -- 7 classes lexicais principais (descr\'edito,
estereotipiza\c{c}\~ao, ass\'edio sexual, amea\c{c}as de viol\^encia,
domina\c{c}\~ao, culpabiliza\c{c}\~ao da v\'itima e objetifica\c{c}\~ao
sexual), derivadas de \cite{anzovino2018} e \cite{sultana2021}, mais 2
subclasses transversais opcionais (\emph{misogynoir} e viol\^encia
pol\'itica digital) -- serve nesta vers\~ao do artigo apenas como
\emph{enquadramento conceitual} para interpreta\c{c}\~ao qualitativa dos
termos induzidos. A avalia\c{c}\~ao reportada aqui \'e exclusivamente
bin\'aria.

\subsection{Corpora e prepara\c{c}\~ao dos dados}

Para a constru\c{c}\~ao do recurso lexical, foram utilizados corpora
anotados em portugu\^es brasileiro previamente consolidados em um corpus
bin\'ario unificado, conforme a Tabela~\ref{tab:corpora}.

\begin{table}[ht]
\centering
\caption{Corpora utilizados na constru\c{c}\~ao do l\'exico, com origem
do r\'otulo de misoginia e quantidade aproveitada ap\'os limpeza.}
\label{tab:corpora}
\begin{tabular}{lllr}
\toprule
\textbf{Corpus} & \textbf{Origem} & \textbf{R\'otulo utilizado} &
\textbf{Docs.} \\
\midrule
HateBR \cite{vargas2022hatebr}      & Instagram & misoginia/\'odio
                                                 hier\'arquico & 7.000 \\
ToLD-BR \cite{leite2020toldbr}      & Twitter   & coluna \texttt{misogyny}
                                                                 & 21.000 \\
Demais fontes agregadas             & misto     & bin\'ario consolidado
                                                                 & 46.338 \\
\midrule
\textbf{Total unificado}            &           &
                                                & \textbf{74.338} \\
\bottomrule
\end{tabular}
\end{table}

O mapeamento dos r\'otulos originais para a distin\c{c}\~ao bin\'aria
adotada seguiu o crit\'erio: documentos com qualquer r\'otulo positivo de
misoginia ou de \'odio direcionado a mulheres foram marcados como
classe~1; os demais como classe~0. Casos amb\'iguos ou textos
excessivamente curtos (\textless\,5 caracteres) foram descartados antes
da consolida\c{c}\~ao final.

\subsection{Pr\'e-processamento e indu\c{c}\~ao do l\'exico}

Os textos passaram por normaliza\c{c}\~ao (caixa baixa, remo\c{c}\~ao de
URLs, men\c{c}\~oes, hashtags e caracteres n\~ao alfab\'eticos),
tokeniza\c{c}\~ao, remo\c{c}\~ao de \emph{stopwords} do portugu\^es e
\textbf{lematiza\c{c}\~ao} pelo spaCy \cite{honnibal2017spacy} com o
modelo \texttt{pt\_core\_news\_lg}. Optou-se por lematiza\c{c}\~ao em vez
de \emph{stemming} porque o recurso ser\'a submetido a valida\c{c}\~ao
manual; lemas s\~ao linguisticamente mais leg\'iveis e reduzem o esfor\c{c}o
cognitivo dos anotadores.

A indu\c{c}\~ao do l\'exico segue o m\'etodo de Monteles
\cite{monteles2023}, em duas fases:

\paragraph{Fase 1: Semente por TF--IDF discriminativo.}
Para cada lema $x$ do vocabul\'ario, calcula-se o TF--IDF separadamente
nas duas subdivis\~oes do corpus (classe~1 vs.\ classe~0). A pontua\c{c}\~ao
discriminativa \'e dada por
$\operatorname{score}(x) = \operatorname{tfidf}_{mis}(x) /
(\operatorname{tfidf}_{non}(x) + \varepsilon)$, e seleciona-se o
\emph{top}-200 por classe -- produzindo 200 sementes para classe~1 e 200
para classe~0.

\paragraph{Fase 2: Expans\~ao por PMI.}
Para cada candidato $x$ do corpus restante, calculam-se as
informa\c{c}\~oes m\'utuas pontuais $\operatorname{PMI}(x, S_{mis})$ e
$\operatorname{PMI}(x, S_{non})$, onde $S_{mis}$ e $S_{non}$ s\~ao os
conjuntos-semente da fase 1. O escore de Monteles \'e

\begin{equation}
O(x) = \operatorname{PMI}(x, S_{mis}) - \operatorname{PMI}(x, S_{non}).
\label{eq:o}
\end{equation}

Aceitam-se candidatos com $|O(x)| \geq 1{,}5$. Os par\^ametros
operacionais foram: janela de coocorr\^encia de 5 lemas, frequ\^encia
m\'inima 5, coocorr\^encia m\'inima 3.

\subsection{Avalia\c{c}\~ao preliminar}

Para avaliar a utilidade do l\'exico, foi adotada uma tarefa bin\'aria
de classifica\c{c}\~ao de textos, comparando-se a configura\c{c}\~ao de
refer\^encia (TF--IDF puro de lemas, \emph{unigramas}, m\'aximo de 10.000
features) com a configura\c{c}\~ao apoiada pelo recurso lexical: \`as
features de TF--IDF adicionam-se tr\^es atributos derivados do l\'exico
calculados por documento -- \texttt{mean\_score} (m\'edia dos escores
$O(x)$ dos termos do documento que aparecem no l\'exico),
\texttt{misog\_ratio} (raz\~ao entre termos com $O(x) > 0$ e o total de
matches) e \texttt{match\_count} (n\'umero de matches no documento).

O conjunto de teste \emph{out-of-sample} \'e composto por 333 frases do
ToLD-BR n\~ao usadas na constru\c{c}\~ao do l\'exico, selecionadas com
crit\'erio de alta confian\c{c}a na anota\c{c}\~ao (\(\geq 2\) votos de
misoginia entre os anotadores originais para a classe positiva, e nenhum
r\'otulo de toxicidade ativo para a classe negativa). Os modelos
comparados s\~ao implementados em scikit-learn \cite{pedregosa2011sklearn}
com \texttt{class\_weight=\textquotesingle balanced\textquotesingle} para
lidar com o desbalanceamento $\sim 80/20$. As m\'etricas reportadas s\~ao
\textbf{F1}, \textbf{precis\~ao} e \textbf{recall} na classe~1 (com
tra\c{c}os de misoginia).

% ===========================================================================
\section{Resultados Preliminares}
% ===========================================================================

A aplica\c{c}\~ao do pipeline produziu um l\'exico de \textbf{531 termos}
em portugu\^es brasileiro -- 268 com escore $O(x) > 0$ (associados \`a
classe~1) e 263 com $O(x) < 0$ (associados \`a classe~0). A
Tabela~\ref{tab:exemplos} ilustra termos representativos do
\emph{top}~20 da classe positiva, com sugest\~ao de classe lexical
obtida por similaridade vetorial \emph{a posteriori}.

\begin{table}[ht]
\centering
\caption{Exemplos selecionados do l\'exico induzido (classe positiva).
Os termos representam \emph{candidatos lexicais discriminativos}: a
classifica\c{c}\~ao final como elemento de um texto com tra\c{c}os de
misoginia depende sempre da ocorr\^encia em contexto.}
\label{tab:exemplos}
\begin{tabular}{lrrl}
\toprule
\textbf{Termo} & \textbf{$O(x)$} & \textbf{$\mathit{freq}_{mis}$} &
\textbf{Classe lexical sugerida} \\
\midrule
canalha      & 4{,}07 & 105 & Descr\'edito \\
feminista    & 3{,}42 &  90 & Descr\'edito \\
bandido      & 3{,}31 & 138 & Amea\c{c}as de Viol\^encia \\
safar        & 3{,}30 &  41 & Amea\c{c}as de Viol\^encia \\
sapat\~ao    & 2{,}98 &  29 & Ass\'edio Sexual \\
piranha      & 2{,}71 &  37 & Ass\'edio Sexual \\
mulherzinha  & 2{,}55 &  17 & Estereotipiza\c{c}\~ao \\
vadia        & 2{,}48 &  44 & Ass\'edio Sexual \\
\bottomrule
\end{tabular}
\end{table}

A Tabela~\ref{tab:resultados} resume os resultados obtidos na
avalia\c{c}\~ao preliminar, comparando a configura\c{c}\~ao baseline
(TF--IDF puro) com a configura\c{c}\~ao apoiada pelo l\'exico, em quatro
classificadores supervisionados.

\begin{table}[ht]
\centering
\caption{Desempenho no corpus de teste \emph{out-of-sample} (333
documentos), classe positiva. F1, precis\~ao e recall referem-se \`a
classe com tra\c{c}os de misoginia. Melhor F1 em \textbf{negrito}.}
\label{tab:resultados}
\begin{tabular}{lccc}
\toprule
\textbf{Modelo} & \textbf{F1} & \textbf{Precis\~ao} & \textbf{Recall} \\
\midrule
SVM linear (TF--IDF + l\'exico)         & \textbf{0,871} & 0,910 & 0,835 \\
Regress\~ao Log\'istica (TF--IDF + l\'exico) & 0,846 & 0,920 & 0,782 \\
Random Forest (TF--IDF + l\'exico)      & 0,560 & 0,718 & 0,459 \\
MLP (TF--IDF + l\'exico)                & 0,506 & 1,000 & 0,338 \\
\bottomrule
\end{tabular}
\end{table}

De modo geral, os resultados indicam que o recurso lexical produzido
apresenta potencial como sinal complementar para a identifica\c{c}\~ao de
textos com tra\c{c}os de misoginia. Os classificadores lineares (SVM e
Regress\~ao Log\'istica) generalizam bem para o conjunto de teste
\emph{out-of-sample}, com F1 acima de 0{,}84. O ganho relativo \`a
baseline -- omitida aqui em favor de uma \'unica tabela principal --
concentra-se sobretudo na \textbf{precis\~ao}, sugerindo que as
\emph{features} derivadas do l\'exico ajudam o classificador a evitar
falsos positivos em textos com vocabul\'ario polissemicamente ambivalente.

\textbf{Limita\c{c}\~oes.} O estudo apresenta as seguintes
limita\c{c}\~oes que delimitam o escopo das conclus\~oes:
(i)~o l\'exico foi induzido a partir de corpora majoritariamente do
Twitter e Instagram, o que limita a generaliza\c{c}\~ao para outros
dom\'inios;
(ii)~muitos termos t\^em forte depend\^encia contextual, refletida na
presen\c{c}a de termos pol\'iticos no \emph{top} do l\'exico que s\'o se
caracterizam como elementos de discurso mis\'ogino quando direcionados a
mulheres em pap\'eis pol\'iticos -- a anota\c{c}\~ao manual em curso pelo
grupo de pesquisa registra esse grau de depend\^encia contextual termo a
termo;
(iii)~a valida\c{c}\~ao humana final do recurso pelos anotadores ainda
n\~ao foi conclu\'ida no momento da reda\c{c}\~ao deste artigo.

% ===========================================================================
\section{Conclus\~ao}
% ===========================================================================

Este artigo apresentou uma abordagem preliminar para constru\c{c}\~ao
semi-autom\'atica de um l\'exico de misoginia em portugu\^es brasileiro
e uma avalia\c{c}\~ao inicial de seu uso em classifica\c{c}\~ao bin\'aria
de textos. O pipeline combina sele\c{c}\~ao por TF--IDF discriminativo,
expans\~ao por PMI e lematiza\c{c}\~ao com spaCy, gerando um recurso
preliminar de 531 termos. A avalia\c{c}\~ao bin\'aria \emph{out-of-sample}
indica desempenho s\'olido (F1 = 0{,}87) com classificadores lineares
quando o l\'exico \'e usado como sinal complementar \`a representa\c{c}\~ao
TF--IDF.

Os resultados sugerem que o recurso lexical constru\'ido \'e promissor
para apoiar tarefas de detec\c{c}\~ao autom\'atica, ainda que dependa de
refinamento metodol\'ogico e valida\c{c}\~ao adicional. Como
continuidade da pesquisa, prev\^em-se: (i)~a consolida\c{c}\~ao da
valida\c{c}\~ao manual do l\'exico pelos anotadores do grupo, com
atribui\c{c}\~ao de classe lexical e grau de depend\^encia contextual
por termo; (ii)~a amplia\c{c}\~ao do corpus de avalia\c{c}\~ao para
textos jornal\'isticos e de comentaristas, fora dos dom\'inios Twitter
e Instagram; e (iii)~a explora\c{c}\~ao de modelos neurais densos quando
houver disponibilidade computacional adequada, para comparar com a
linha de base linear apresentada aqui.

\bibliographystyle{sbc}
\bibliography{referencias}

\end{document}
